\ifdefined\gkver\else\def\gkver{student}\fi
\documentclass[\gkver]{gaokaozhenti}
\begin{document}

\chapter{2026年天津}

\begin{enumerate}
\item
%% number: 1
%% paperName: 2026年天津高考第 1 题 5 分
%% typeId: 1
%% type: 单选题
%% chapter: 原子和原子核
%% point: 人工核反应
%% method: 无
%% score: 5
%% degree: 850
%% duplicateId: 0
%% body:
2026 年中国超重元素研究加速器装置取得了新突破。利用该装置加速的 $\ce{^{40}_{18}Ar}$ 原子核轰击 $\ce{^{197}_{79}Au}$ 原子核，产生元素 X 和 2 个中子，其核反应方程为 $\ce{^{40}_{18}Ar + ^{197}_{79}Au -> X + 2^{1}_{0}n}$，则\xzanswer{C}
\fourchoices[answer=C]
{该方程是衰变方程}
{该反应为核裂变反应}
{X 原子核的质子数为 97}
{X 原子核的中子数为 235}
%% answer: C
%% memo:
\memoanswer{
A．衰变的本质是原子核自发地放出 $\alpha$、$\beta$ 等粒子的核反应过程，该反应是 Ar 原子核轰击 Au 原子核诱发的人工核反应，并非原子核自发发生的衰变，故 A 错误；

B．核裂变是重核分裂为多个中等质量原子核的反应，该反应是两个原子核结合生成更重的新核的人工核反应，不符合裂变的定义，故 B 错误；

C．根据核反应的电荷数守恒规律，反应前总电荷数为 $18 + 79 = 97$，反应后 2 个中子的电荷数均为 0，因此 X 原子核的质子数等于总电荷数 97，故 C 正确；

D．根据核反应的质量数守恒规律，反应前总质量数为 $40 + 197 = 237$，反应后 2 个中子的总质量数为 $2\times1 = 2$，因此 X 原子核的质量数为 $237 - 2 = 235$；原子核的中子数为质量数与质子数的差值，即 $235 - 97 = 138$，故 D 错误。

故选 C。
}

\item
%% number: 2
%% paperName: 2026年天津高考第 2 题 5 分
%% typeId: 1
%% type: 单选题
%% chapter: 气体性质
%% point: 热力学第一定律
%% method: 无
%% score: 5
%% degree: 750
%% duplicateId: 0
%% body:
公交车利用高压气体推动活塞实现车门开启。如图所示，打开阀门后，储气罐内的高压气体进入气缸活塞左侧的密封区域，推动活塞向右运动，此过程视为绝热过程。高压气体视为理想气体，则气体\xzanswer{B}
\onepicture[w=7.5cm]{\includesvg{figs/02}}
\fourchoices[answer=B]
{温度不变}
{内能减少}
{压强不变}
{对外不做功}
%% answer: B
%% memo:
\memoanswer{
D．推动活塞向右运动，可知气体体积变大，则气体对外做功，故 D 错误；

AB．由于气体体积变大，外界对气体做负功，且此过程视为绝热过程，根据热力学第一定律可知，气体内能减少，则气体温度降低，故 A 错误，B 正确；

C．根据理想气体状态方程 $\frac{pV}{T} = C$，由于气体体积变大，温度降低，所以气体压强减小，故 C 错误。

故选 B。
}

\item
%% number: 3
%% paperName: 2026年天津高考第 3 题 5 分
%% typeId: 1
%% type: 单选题
%% chapter: 光学
%% point: 光电效应
%% method: 无
%% score: 5
%% degree: 650
%% duplicateId: 0
%% body:
如图所示，一束复色光从空气射入玻璃三棱镜后，出射光分成 a、b 两束单色光，a 的波长大于 b 的波长，则\xzanswer{D}
\onepicture[w=5cm]{\includesvg{figs/03}}
\fourchoices[answer=D]
{a 的光子动量大于 b 的}
{在三棱镜中，a 的传播速度小于 b 的}
{a、b 在空间互相叠加，可发生干涉现象}
{照射同一金属，若 a 能发生光电效应，则 b 也一定能}
%% answer: D
%% memo:
\memoanswer{
A．光子动量公式为 $p = \frac{h}{\lambda}$，波长越大动量越小，因此 a 的光子动量小于 b 的，故 A 错误；

B．由图可知 a 光的偏折程度较小，则折射率较小，根据 $v = \frac{c}{n}$ 可知，a 在三棱镜中传播速度较大，故 B 错误；

C．光发生干涉的必要条件是两束光频率相同，而两束光折射率不同，则频率不同，则不能发生干涉，故 C 错误；

D．由上分析可知，a 光折射率较小，频率较小，若其能使金属发生光电效应，说明其频率大于该金属的截止频率，则频率更大的 b 光肯定也能使该金属发生光电效应，故 D 正确。

故选 D。
}

\item
%% number: 4
%% paperName: 2026年天津高考第 4 题 5 分
%% typeId: 1
%% type: 单选题
%% chapter: 机械波
%% point: 波与振动图象
%% method: 无
%% score: 5
%% degree: 650
%% duplicateId: 0
%% body:
一列简谐横波在均匀介质中沿 $x$ 轴传播，$t = 0$ 时的波形如图所示。介质中 $x = 2 \Um$ 处的质点 $P$ 做简谐运动的表达式为 $y = 5\sin(2\pi t) \Ucm$。则\xzanswer{D}
\onepicture[w=8cm]{\includesvg{figs/04}}
\fourchoices[answer=D]
{该波的周期为 $2 \Us$}
{该波沿 $x$ 轴的负方向传播}
{该波在此介质中的波速为 $2 \Ums$}
{该时刻质点 $P$ 沿 $y$ 轴正方向运动}
%% answer: D
%% memo:
\memoanswer{
AC．由 $y = 5\sin(2\pi t) \Ucm$，可知该波的周期为 $T = \frac{2\pi}{\omega} = \frac{2\pi}{2\pi} \Us = 1 \Us$。

由题图可知波长为 $\lambda = 4 \Um$，则波速为 $v = \frac{\lambda}{T} = 4 \Ums$，故 AC 错误；

BD．已知质点 $P$ 做简谐运动的表达式为 $y = 5\sin(2\pi t) \Ucm$，可知 $t = 0$ 时刻质点 $P$ 沿 $y$ 轴正方向运动，根据波形平移法可知，该波沿 $x$ 轴的正方向传播，故 B 错误，D 正确。

故选 D。
}

\item
%% number: 5
%% paperName: 2026年天津高考第 5 题 5 分
%% typeId: 1
%% type: 单选题
%% chapter: 曲线运动
%% point: 线速度角速度
%% method: 无
%% score: 5
%% degree: 650
%% duplicateId: 0
%% body:
隋唐时期出现的井车是一种从深井提水灌溉的农具。示意如图，井车主要由齿轮 a、齿轮 b 和若干个水斗构成。工作时，由 a 匀速转动带动 b 转动提升水斗，将水提升至地面灌溉农田。a 与 b 的半径之比为 $r_{a}:r_{b}$，在井车工作过程中，下列说法正确的是\xzanswer{A}
\onepicture[w=5.5cm]{figs/05.png}
\fourchoices[answer=A]
{a、b 角速度的大小之比为 $r_{b}:r_{a}$}
{a、b 边缘线速度大小之比为 $r_{a}:r_{b}$}
{图中最低处水斗所受合外力为零}
{水斗运动一周所受重力的冲量为零}
%% answer: A
%% memo:
\memoanswer{
AB．齿轮 a、b 属于啮合传动，两齿轮边缘的线速度大小相等，即 $v_{a} = v_{b}$。

根据 $v = \omega r$ 可知，a、b 角速度的大小之比为 $\omega_{a}:\omega_{b} = r_{b}:r_{a}$，故 A 正确，B 错误；

C．水斗随齿轮 b 做匀速圆周运动，合外力提供向心力，指向圆周运动的圆心，最低处水斗的合外力竖直向上、不为零，故 C 错误；

D．根据冲量的定义可得 $I_{G} = mgt$，可知水斗运动一周所受重力的冲量不为零，故 D 错误。

故选 A。
}

\item
%% number: 6
%% paperName: 2026年天津高考第 6 题 5 分
%% typeId: 2
%% type: 多选题
%% chapter: 曲线运动
%% point: 人造地球卫星
%% method: 无
%% score: 5
%% degree: 550
%% duplicateId: 0
%% body:
探测器绕行星运动视为匀速圆周运动，探测器的质量为 $m$、轨道半径为 $r$、运行周期为 $T$，引力常量为 $G$，不考虑行星自转的影响，则\xzanswer{AB}
\fourchoices[answer=AB]
{行星质量 $M = \frac{4\pi^{2}r^{3}}{GT^{2}}$}
{探测器线速度 $v = \frac{2\pi r}{T}$}
{行星表面自由落体加速度 $g = \frac{4\pi^{2}r^{2}}{T^{2}}$}
{探测器向心加速度 $a = \frac{Gm}{r^{2}}$}
%% answer: AB
%% memo:
\memoanswer{
A．根据 $\frac{GMm}{r^{2}} = m\frac{4\pi^{2}}{T^{2}}r$，可得行星质量 $M = \frac{4\pi^{2}r^{3}}{GT^{2}}$，A 正确；

B．探测器线速度 $v = \frac{2\pi r}{T}$，B 正确；

C．设行星的半径为 $R$，根据 $\frac{GMm_{0}}{R^{2}} = m_{0}g$，可得行星表面自由落体加速度 $g = \frac{4\pi^{2}r^{3}}{R^{2}T^{2}}$，C 错误；

D．根据 $\frac{GMm}{r^{2}} = ma$，可得探测器向心加速度 $a = \frac{GM}{r^{2}} = \frac{4\pi^{2}r}{T^{2}}$，D 错误。

故选 AB。
}

\item
%% number: 7
%% paperName: 2026年天津高考第 7 题 5 分
%% typeId: 2
%% type: 多选题
%% chapter: 交流电
%% point: 变压器
%% method: 无
%% score: 5
%% degree: 450
%% duplicateId: 0
%% body:
在生产实践中，需要保证用电器电压稳定。如图所示，理想变压器的原线圈接在电压有效值不变的正弦交流电源上，匝数保持不变。副线圈两端连接了定值电阻 $R$ 和若干并联的用电器，当用电器数量发生变化时，可自动调节滑动触头 Q 的位置改变副线圈匝数，保持用电器两端电压不变。则用电器数量增加时\xzanswer{BC}
\onepicture[w=6.5cm]{\includesvg{figs/07}}
\fourchoices[answer=BC]
{Q 向下滑动}
{$R$ 两端电压变大}
{原线圈的电流变大}
{电源输出功率不变}
%% answer: BC
%% memo:
\memoanswer{
B．用电器数量增加时，则用电器总电阻 $R_{\text{用}}$ 减小，已知用电器两端电压不变，根据欧姆定律可得 $I_{2} = \frac{U_{\text{用}}}{R_{\text{用}}}$，可知副线圈电流 $I_{2}$ 增大，即通过 $R$ 的电流增大，所以 $R$ 两端电压变大，故 B 正确；

A．根据 $U_{2} = I_{2}R + U_{\text{用}}$，由于副线圈电流 $I_{2}$ 增大，则副线圈两端电压 $U_{2}$ 增大；根据 $\frac{U_{2}}{U_{1}} = \frac{n_{2}}{n_{1}}$ 可知，副线圈匝数增大，则 Q 向上滑动，故 A 错误；

CD．根据 $P = U_{1}I_{1} = U_{2}I_{2}$，由于 $U_{2}$、$I_{2}$ 增大，则电源输出功率增大；由于原线圈输入电压 $U_{1}$ 不变，所以原线圈的电流 $I_{1}$ 变大，故 C 正确，D 错误。

故选 BC。
}

\item
%% number: 8
%% paperName: 2026年天津高考第 8 题 5 分
%% typeId: 2
%% type: 多选题
%% chapter: 电场
%% point: 带电粒子的直线运动
%% method: 无
%% score: 5
%% degree: 450
%% duplicateId: 0
%% body:
如图所示，$O$ 点为均匀带负电固定圆环的圆心，由圆环中心轴线上某点 $A$ 静止释放一带正电的试探电荷 $q$，$q$ 只在静电力作用下运动，则\xzanswer{AC}
\onepicture[w=3.6cm]{\includesvg{figs/08}}
\fourchoices[answer=AC]
{$q$ 在 $O$ 点所受静电力为零}
{取无限远处电势为零，$O$ 点电势也为零}
{$q$ 由 $A$ 向 $O$ 运动过程中，动能一定逐渐增大}
{$q$ 由 $A$ 向 $O$ 运动过程中，加速度一定逐渐减小}
%% answer: AC
%% memo:
\memoanswer{
A．均匀带负电圆环的圆心 $O$ 处，圆环上每个电荷元对试探电荷的静电力，都会关于 $O$ 点对称抵消，合场强为 0，因此带正电的试探电荷在 $O$ 点受到的静电力为零，故 A 正确；

B．取无限远电势为零时，将一个正试探电荷从 $O$ 点移动到无限远，带负电的圆环对正试探电荷的静电力做负功，电势能增加，说明 $O$ 点电势低于无限远的电势，则 $O$ 点电势为负值，并非为零，故 B 错误；

C．带负电圆环的中心轴线上，$O$ 点场强为 0，轴线上 $O$ 点外侧（$A$ 所在的一侧）的场强方向是指向 $O$ 点。带正电的试探电荷从 $A$ 向 $O$ 运动时，静电力方向与运动方向始终一致，静电力做正功，动能一定逐渐增大，故 C 正确；

D．轴线上从 $O$ 点向无穷远延伸时，场强从 0 开始，先增大到一个最大值，再逐渐减小到 0。如果 $A$ 点位于场强最大值的外侧，从 $A$ 向 $O$ 运动的过程中，场强会先增大、后减小，由 $a = \frac{qE}{m}$ 可知，加速度也会先增大后减小，并非一定逐渐减小，故 D 错误。

故选 AC。
}

\item
%% number: 9
%% paperName: 2026年天津高考第 9 题 6 分
%% typeId: 4
%% type: 实验题
%% chapter: 功和能
%% point: DIS研究机械能守恒定律
%% method: 无
%% score: 6
%% degree: 650
%% duplicateId: 0
%% body:
利用如图~\ref{2026天津09a}所示装置验证机械能守恒定律。
\begin{enumerate}
	\item
	打点计时器固定在铁架台上，使重物带动纸带从静止开始自由下落。某次实验得到一条纸带，部分计数点如图~\ref{2026天津09b}所示，相邻点间的时间间隔为 $t$，测得 $B$、$C$ 点间距离为 $x_{1}$，$C$、$D$ 点间距离为 $x_{2}$，则 $C$ 点的速度大小为\tkanswer{$\frac{x_{1} + x_{2}}{2t}$}，图中所示的纸带\tkanswer{左端}与重物相连（选填“左端”或“右端”）；
	\item
	将点 $E$ 对应重物所在的高度取为 0，测量纸带上各点到点 $E$ 间的距离，由此获得各点对应重物所在的高度 $h$ 和速度大小 $v$，绘制出 $\frac{v^{2}}{2}-h$ 图像如图~\ref{2026天津09c}所示，在误差允许范围内视为机械能守恒，则图中直线斜率表示的物理量是\tkanswer{$-g$}。
\end{enumerate}
\threepicture[labelA=2026天津09a, labelB=2026天津09b, labelC=2026天津09c, wA=3.0cm, wB=8.0cm, wC=3.2cm, align=right]{\includesvg{figs/09a}}{\includesvg{figs/09b}}{\includesvg{figs/09c}}
%% answer:
\jdanswer{
\begin{enumerate}
	\item
	$\frac{x_{1} + x_{2}}{2t}$，左端
	\item
	$-g$
\end{enumerate}
}
%% memo:
\memoanswer{
（1）$C$ 点的速度大小为 $v_{C} = \frac{x_{1} + x_{2}}{2t}$，因纸带上从左到右相邻点间距逐渐增加，可知纸带左端与重物相连。

（2）将点 $E$ 对应重物所在的高度取为 0，各点对应重物所在的高度 $h$，设 $E$ 点时重物下落的速度为 $v_{0}$，则若机械能守恒，则根据 $mgh = \frac{1}{2}mv_{0}^{2} - \frac{1}{2}mv^{2}$，可得 $\frac{v^{2}}{2} = \frac{v_{0}^{2}}{2} - gh$，则图中直线斜率表示的物理量是 $-g$。
}

\item
%% number: 9
%% paperName: 2026年天津高考第 9 题 6 分
%% typeId: 4
%% type: 实验题
%% chapter: 电场
%% point: 电容器充放电实验
%% method: 无
%% score: 6
%% degree: 550
%% duplicateId: 0
%% body:
实验小组借助电流和电压传感器及计算机等设备，利用如图~\ref{2026天津09d}所示电路观察电容器的充、放电过程。
\begin{enumerate}
	\item
	下列能正确描述电容器电荷量 $Q$、电压 $U$、电容 $C$ 间的关系图像是\tkanswer{B}；
	\item
	实验中，先使开关 S 与 1 端相连，电源向电容器充电，然后把开关 S 掷向 2 端，电容器放电，传感器将电流 $I$ 和电压 $U$ 随时间 $t$ 变化的数据传入计算机，屏幕上显示出充电和放电全过程的 $I-t$、$U-t$ 图像如图~\ref{2026天津09f}所示（图像纵轴未标出），则\tkanswer{B}是 $I-t$ 图像（选填“A”或“B”）；
	\item
	选取电容为 $C_{1}$、$C_{2}$ 的两个电容器，分别接入图~\ref{2026天津09d}所示电路，得到放电过程的电流大小随时间变化的关系图像如图~\ref{2026天津09g}中的实线、虚线所示，则 $C_{1}$\tkanswer{小于}$C_{2}$（选填“大于”、“小于”或“等于”）。
\end{enumerate}
\fourpicture[labelA=2026天津09d, labelB=2026天津09e, labelC=2026天津09f, labelD=2026天津09g, wA=6.0cm, wB=8.0cm, wC=9.8cm, wD=4.0cm, align=right]{\includesvg{figs/09d}}{\includesvg{figs/09e}}{\includesvg{figs/09f}}{\includesvg{figs/09g}}
%% answer:
\jdanswer{
\begin{enumerate}
	\item
	B
	\item
	B
	\item
	小于
\end{enumerate}
}
%% memo:
\memoanswer{
（1）AC．电容器的电容 $C$ 与电容器所带的电量 $Q$ 以及两极板的电势差 $U$ 无关，则 $C-Q$ 图像是平行于 $Q$ 轴的直线，$C-U$ 图像是平行于 $U$ 轴的直线，AC 错误；

B．根据 $Q = CU$ 可知 $Q-U$ 图像为过原点的倾斜直线，B 正确。

故选 B。

（2）电容器充电时电流从最大值逐渐减小到零，放电时电流反向且从最大值逐渐减小到零，所以通过传感器的电流 $I$ 随时间 $t$ 变化的图像为 B。

（3）因图像与坐标轴围成的面积等于电容器的带电量，因虚线与坐标轴围成的面积较大，可知 $C_{2}$ 电容器带电量较大，根据 $Q = CU$ 可知，$C_{1}$ 小于 $C_{2}$。
}

\item
%% number: 10
%% paperName: 2026年天津高考第 10 题 14 分
%% typeId: 6
%% type: 计算题
%% chapter: 功和能
%% point: 动能定理
%% method: 无
%% score: 14
%% degree: 550
%% duplicateId: 0
%% body:
如图所示，一根长 $l = 0.4 \Um$ 的轻绳 1 一端固定，另一端系着一质量 $m = 0.02 \Ukg$ 的小球。初始时，小球在绳 1 及水平轻绳 2 共同作用下保持静止状态，绳 1 与竖直方向夹角 $\theta = 60^{\circ}$。现剪断绳 2，小球无初速度下摆至悬挂点的正下方，与静置在水平地面上的木块发生碰撞，碰撞时间极短，碰后小球速度为零。木块质量 $M = 0.04 \Ukg$，与地面的动摩擦因数 $\mu = 0.1$。小球和木块均视为质点，轻绳不可伸长，$g$ 取 $10 \Umsq$，忽略空气阻力。求
\begin{enumerate}
	\item
	剪断绳 2 前，绳 1 对小球拉力 $F$ 的大小；
	\item
	小球与木块碰撞前瞬间，小球速度 $v$ 的大小；
	\item
	木块在地面上运动位移 $x$ 的大小。
\end{enumerate}
\onepicture[w=6.5cm, align=right]{\includesvg{figs/10}}
%% answer:
\jdanswer{
\begin{enumerate}
	\item
	$F = 0.4 \UN$
	\item
	$v = 2 \Ums$
	\item
	$x = 0.5 \Um$
\end{enumerate}
}
%% memo:
\memoanswer{
（1）剪断绳 2 前，以小球为对象，根据平衡条件可得 $\cos\theta = \frac{mg}{F}$，可得绳 1 对小球的拉力大小为 $F = \frac{mg}{\cos\theta} = 0.4 \UN$。

（2）剪断绳 2，小球向下摆动过程中，根据动能定理可得 $mgl(1 - \cos\theta) = \frac{1}{2}mv^{2} - 0$，解得小球与木块碰撞前瞬间，小球的速度大小为 $v = \sqrt{2gl(1 - \cos\theta)} = 2 \Ums$。

（3）小球与木块碰撞过程，根据动量守恒可得 $mv = Mv_{1}$，解得碰后木块的速度为 $v_{1} = 1 \Ums$。

对木块，根据动能定理可得 $-\mu Mgx = 0 - \frac{1}{2}Mv_{1}^{2}$，解得木块在地面上运动的位移大小为 $x = \frac{v_{1}^{2}}{2\mu g} = 0.5 \Um$。
}

\item
%% number: 11
%% paperName: 2026年天津高考第 11 题 16 分
%% typeId: 6
%% type: 计算题
%% chapter: 电磁感应
%% point: 电磁感应的电量
%% method: 无
%% score: 16
%% degree: 450
%% duplicateId: 0
%% body:
如图所示，将两个完全相同的正方形金属框的上、下边均弯曲成半径为 $r$ 的半圆，半圆所在平面与金属框左、右两边垂直，两框分别记为 $abcd$ 和 $a'b'c'd'$，两框正对放置并用绝缘材料连接成一个整体，彼此绝缘。金属框下方有高为 $h$（$h < \pi r$）且足够长和宽的水平匀强磁场区域，磁感应强度大小为 $B$，方向垂直于 $ab$、$cd$ 构成的平面。已知金属框整体以竖直向下的速度 $v$ 进入磁场，$abcd$ 和 $a'b'c'd'$ 电阻均为 $R$，下落过程中金属框整体不发生转动，不计空气阻力。
\begin{enumerate}
	\item
	判断金属框刚进入磁场时 $ab$ 边的电流方向；
	\item
	求金属框 $bc$ 边穿过磁场区域的过程中，通过 $bc$ 边横截面的电荷量 $q$；
	\item
	求金属框刚进入磁场时整体受到的安培力大小 $F$。
\end{enumerate}
\onepicture[w=10cm, align=right]{\includesvg{figs/11}}
%% answer:
\jdanswer{
\begin{enumerate}
	\item
	电流方向由 $a$ 流向 $b$
	\item
	$q = \frac{2Brh}{R}$
	\item
	$F = \frac{8B^{2}r^{2}v}{R}$
\end{enumerate}
}
%% memo:
\memoanswer{
（1）金属框刚进入磁场时，金属框 $bc$ 边在磁场中向下切割磁感线，根据右手定则可知，金属框的电流方向为 $bcdab$，则 $ab$ 边的电流方向为由 $a$ 流向 $b$。

（2）根据 $q = \bar{I}\Delta t = \frac{\Delta\Phi}{R} = \frac{B\Delta S}{R}$，其中 $\Delta S$ 为垂直磁场方向扫过的面积，则金属框 $bc$ 边穿过磁场区域的过程中，有 $\Delta S = 2rh$，联立可得通过 $bc$ 边横截面的电荷量为 $q = \frac{2Brh}{R}$。

（3）金属框刚进入磁场时，金属框 $bc$ 边、$b'c'$ 边切割磁感线，有效长度均为 $l = 2r$，则 $bc$ 边、$b'c'$ 边产生的电动势均为 $E = Blv = 2Brv$，两个金属框的电流均为 $I = \frac{E}{R} = \frac{2Brv}{R}$。

$bc$ 边、$b'c'$ 边受到的安培力大小均为 $F_{1} = BIl = \frac{4B^{2}r^{2}v}{R}$，则金属框刚进入磁场时整体受到的安培力大小为 $F = 2F_{1} = \frac{8B^{2}r^{2}v}{R}$。
}

\item
%% number: 12
%% paperName: 2026年天津高考第 12 题 18 分
%% typeId: 6
%% type: 计算题
%% chapter: 磁场
%% point: 带电粒子在磁场中的运动
%% method: 无
%% score: 18
%% degree: 350
%% duplicateId: 0
%% body:
磁透镜利用磁场对带电粒子进行聚焦、成像，是电子显微镜等科学装置的核心部件。某磁透镜圆柱形磁场区域如图 1 所示，长度为 $L$，对称轴为 $z$ 轴，区域外磁场忽略不计；图 2 为图 1 过 $z$ 轴的截面图，坐标原点 $O$ 为圆柱中心，磁场为轴对称磁场，磁感应强度沿 $z$ 轴正方向分量为 $B_{z}$，垂直于 $z$ 轴分量方向以过原点 $O$ 与 $z$ 轴垂直的平面为界，左侧指向 $z$ 轴、右侧背离 $z$ 轴。一初速度忽略不计的电子，质量为 $m$、电荷量为 $e$，经电压为 $U$ 的电场加速后沿 $z$ 轴正方向进入磁场。为研究方便，将电子进入磁场后的运动分解为沿 $z$ 方向的运动及与 $z$ 轴垂直的平面内的运动，相应的速度分解为 $v_{z}$、$v_{r}$、$v_{\varphi}$ 三个分量，其中 $v_{z}$ 沿 $z$ 轴正方向，$v_{r}$ 与 $v_{\varphi}$ 处于与 $z$ 轴垂直的平面内，$v_{r}$ 指向 $z$ 轴，$v_{\varphi}$ 垂直于 $v_{r}$，如图 3 所示（垂直纸面向外为 $z$ 轴正方向）。不计相对论效应、电磁辐射及电子重力。
\begin{enumerate}
	\item
	求电子刚进入磁场时的速度大小 $v_{0}$；
	\item
	计算 $v_{r}$ 时，若将电子在垂直于 $z$ 轴平面内的运动等效为指向 $z$ 轴的直线运动，则电子受到的在该直线上的合力为洛伦兹力在此方向分量的一半。已知 $v_{\varphi} = \frac{re}{2m}B_{z}$（$r$ 为电子到 $z$ 轴距离），$B_{z}^{2}$ 与 $z$ 的关系如图 4 所示，曲线与 $z$ 轴所围面积为 $S$。
	\begin{enumerate}
		\item
		指向 $z$ 轴方向的洛伦兹力分量与哪个速度分量有关；
		\item
		由于 $v_{0}$ 很大，$v_{z} \approx v_{0}$，电子在磁场中运动时间极短、与 $z$ 轴的距离视为不变，求该电子从离开磁场到运动至 $z$ 轴时沿 $z$ 方向的分位移大小 $f$；
	\end{enumerate}
	\item
	通过计算结果分析说明该磁透镜能够对与 $z$ 轴平行入射的电子束聚焦。
\end{enumerate}
\onepicture[w=13cm, align=right]{\includesvg{figs/12}}
%% answer:
\jdanswer{
\begin{enumerate}
	\item
	$v_{0} = \sqrt{\frac{2eU}{m}}$
	\item
	（i）$v_{\varphi}$；（ii）$f = \frac{8mU}{eS}$
	\item
	电子离开磁场时获得的指向 $z$ 轴的径向速度为 $v_{r} = \frac{e^{2}rS}{4m^{2}v_{0}}$。可见 $v_{r}$ 与离轴距离 $r$ 成正比，因此电子从离开磁场到运动到 $z$ 轴所需时间为 $t = \frac{r}{v_{r}} = \frac{4m^{2}v_{0}}{e^{2}S}$，该时间与 $r$ 无关。所以不同离轴距离的平行入射电子，在离开磁场后经过相同的轴向距离 $f = v_{0}t = \frac{4m^{2}v_{0}^{2}}{e^{2}S} = \frac{8mU}{eS}$，之后同时到达 $z$ 轴上的同一点，因此该磁场区域对平行于 $z$ 轴入射的电子束具有聚焦作用，相当于一个磁透镜。
\end{enumerate}
}
%% memo:
\memoanswer{
（1）电子经电压 $U$ 加速，电场力做功等于动能增加，有 $eU = \frac{1}{2}mv_{0}^{2}$，因此 $v_{0} = \sqrt{\frac{2eU}{m}}$，方向沿 $+z$ 方向。

（2）（i）由题意可知，在坐标系中，轴对称磁场无角向分量，即 $B_{\varphi} = 0$，洛伦兹力的径向分量为 $F_{r} = ev_{\varphi}B_{z}$，因此指向 $z$ 轴方向的洛伦兹力分量与速度分量 $v_{\varphi}$ 有关。

（ii）根据题意有 $v_{\varphi} = \frac{re}{2m}B_{z}$，其中 $r$ 为电子到 $z$ 轴的距离。径向洛伦兹力的大小为 $F_{r} = ev_{\varphi}B_{z} = e\cdot\frac{re}{2m}B_{z}^{2} = \frac{e^{2}r}{2m}B_{z}^{2}$。

按题意，取等效径向合力为其一半，即 $F_{\text{等效}} = \frac{1}{2}F_{r} = \frac{e^{2}r}{4m}B_{z}^{2}$，因此径向加速度为 $\frac{\mathrm{d}v_{r}}{\mathrm{d}t} = \frac{F_{\text{等效}}}{m} = \frac{e^{2}r}{4m^{2}}B_{z}^{2}$。

由于 $v_{0}$ 很大，近似有 $v_{z} \approx v_{0}$，且电子在磁场中径向距离 $r$ 近似不变，$\mathrm{d}t = \frac{\mathrm{d}z}{v_{z}} = \frac{\mathrm{d}z}{v_{0}}$，于是电子离开磁场时的径向速度为 $v_{r} = \frac{e^{2}r}{4m^{2}v_{0}}\int B_{z}^{2}\,\mathrm{d}z$。

由图 4，曲线与 $z$ 轴所围面积为 $S = \int B_{z}^{2}\,\mathrm{d}z$，故 $v_{r} = \frac{e^{2}rS}{4m^{2}v_{0}}$，该速度指向 $z$ 轴。电子离开磁场后不再受磁场力，做匀速直线运动，它从离轴距离 $r$ 处运动到 $z$ 轴所需时间为 $t = \frac{r}{v_{r}} = \frac{4m^{2}v_{0}}{e^{2}S}$。

在这段时间内沿 $z$ 方向的位移为 $f = v_{z}t = v_{0}t = \frac{4m^{2}v_{0}^{2}}{e^{2}S} = \frac{8mU}{eS}$。

（3）电子离开磁场时获得的指向 $z$ 轴的径向速度为 $v_{r} = \frac{e^{2}rS}{4m^{2}v_{0}}$，可见 $v_{r}$ 与离轴距离 $r$ 成正比，因此电子从离开磁场到运动到 $z$ 轴所需时间为 $t = \frac{r}{v_{r}} = \frac{4m^{2}v_{0}}{e^{2}S}$，该时间与 $r$ 无关。

所以，不同离轴距离的平行入射电子，在离开磁场后经过相同的轴向距离 $f = v_{0}t = \frac{4m^{2}v_{0}^{2}}{e^{2}S} = \frac{8mU}{eS}$，之后同时到达 $z$ 轴上的同一点。因此该磁场区域对平行于 $z$ 轴入射的电子束具有聚焦作用，相当于一个磁透镜。
}

\end{enumerate}

\end{document}
